Para confirmar que a eq. (5.23) é solução da eq. (5.22), substitui-se diretamente $V_C(t)$ e sua derivada na equação diferencial.

\textbf{Solução proposta (eq. 5.23):}
\[
V_C(t) = V_s\left(1 - e^{-t/(RC)}\right).
\]

\textbf{Cálculo da derivada:}
\[
\frac{dV_C}{dt} = V_s\cdot\frac{1}{RC}\,e^{-t/(RC)}.
\]

\textbf{Substituição na eq. (5.22):}
\begin{align*}
RC\,\frac{dV_C}{dt} + V_C
&= RC\cdot\frac{V_s}{RC}\,e^{-t/(RC)} + V_s\left(1 - e^{-t/(RC)}\right) \\
&= V_s\,e^{-t/(RC)} + V_s - V_s\,e^{-t/(RC)} \\
&= V_s.\quad\checkmark
\end{align*}

\textbf{Condição inicial:} Em $t=0$:
\[
V_C(0) = V_s\left(1 - e^0\right) = V_s(1-1) = 0.\quad\checkmark
\]

\textbf{Comportamento assintótico:} Quando $t \to \infty$, $e^{-t/(RC)} \to 0$ e $V_C \to V_s$, confirmando que o capacitor carrega até a tensão da fonte.

\textbf{Conclusão:} A eq. (5.23) satisfaz identicamente a equação diferencial e a condição inicial, sendo a solução correta. 